\section{伦理与社会影响}

\subsection{AI 的双刃剑效应}

Fei-Fei Li 在演讲结尾强调：强大的工具伴随着巨大的责任。计算机视觉可以造福社会，也可能造成伤害。

\begin{figure}[H]
\centering
\includegraphics[width=0.95\textwidth]{slides-images/slide-065.jpg}
\caption{计算机视觉的潜在危害：左侧为面部识别算法中的性别和肤色偏见（不同系统对深色皮肤女性的识别准确率明显偏低），右侧为 AI 决策对求职者命运的影响\protect\footnotemark}
\end{figure}
\footnotetext{Slides 第66页（Part 1）。}

\begin{warningbox}{数据偏见与 AI 公平性}
当前所有大型 AI 算法都由数据驱动，而数据是人类活动和历史的产物，天然携带了各种\textbf{偏见}。面部识别算法已被多次发现存在种族和性别偏见——对深色皮肤女性的识别错误率远高于浅色皮肤男性。当 AI 被用于影响人们的生活（招聘、贷款审批等）时，这些偏见可能被系统性地放大。
\end{warningbox}

\begin{knowledgebox}{AI 应用的积极一面}
Fei-Fei Li 特别提到了 AI 在医疗健康领域的应用。她和 Ehsan Adeli、Zane Durante 共同研究使用计算机视觉为老龄人口和患者提供照护支持。AI 在医学影像分析等领域可以极大地提升效率和准确性，惠及更多患者。
\end{knowledgebox}

\subsection{人类视觉的启示}

Fei-Fei Li 最后指出，尽管计算机视觉已经取得了巨大进步，但人类视觉的丰富性、微妙性和情感深度仍然远超当前 AI 的能力。理解并欣赏人类视觉的卓越之处，应该是持续激励研究者前进的动力。

\subsection{本章小结}

技术进步必须伴随伦理反思。数据中的人类偏见会被 AI 系统继承和放大，这要求研究者和从业者高度关注公平性、透明度和社会影响。同时，AI 也在医疗等领域展现了巨大的正面潜力。

%% ========== 课程概览 ==========

